\documentclass[a4paper,11pt]{article}

% ============================================================
%  Le Collé — Séance 1 : Déterminants, courbes paramétrées
%  © 2026 Ahmed Amine Ouqlani
% ============================================================

% ---------- Polices ----------
\usepackage[utf8]{inputenc}
\usepackage[T1]{fontenc}
\usepackage[sc]{mathpazo}          % texte et mathématiques : Palatino
\usepackage[scaled=0.88]{helvet}   % titres et repères : linéale
\linespread{1.07}
\usepackage{amsmath,amssymb,amsthm,mathrsfs,stmaryrd}
\usepackage{microtype}

% ---------- Mise en page ----------
\usepackage[a4paper,top=2.3cm,bottom=2.3cm,left=2.4cm,right=2.4cm,
            headheight=14pt,headsep=0.8cm,footskip=1.1cm]{geometry}
\setlength{\parindent}{0pt}
\setlength{\parskip}{0.45em plus 0.15em minus 0.1em}
\raggedbottom
\emergencystretch=1.5em
\clubpenalty=10000 \widowpenalty=10000 \displaywidowpenalty=10000 \brokenpenalty=10000

% ---------- Couleurs : un bleu de nuit, un ocre, des gris ----------
\usepackage{xcolor}
\definecolor{Ink}{HTML}{1C2533}
\definecolor{Bleu}{HTML}{1F3A5F}
\definecolor{BleuMoyen}{HTML}{4A6A8F}
\definecolor{BleuClair}{HTML}{9DB2C9}
\definecolor{Fond}{HTML}{F2F4F7}
\definecolor{Ocre}{HTML}{B07A26}
\definecolor{Brique}{HTML}{A23B2A}
\definecolor{Gris}{HTML}{5E6672}
\definecolor{Filet}{HTML}{C5CCD6}

% ---------- Packages ----------
\usepackage{enumitem}
\setlist{itemsep=0.25em,topsep=0.35em,parsep=0pt}
\setlist[itemize]{label=\textcolor{BleuMoyen}{\small\textbullet},leftmargin=1.4em}
\setlist[enumerate]{leftmargin=1.8em,label=\textsf{\small\color{Bleu}\arabic*.}}
\usepackage{booktabs,array,multicol}
\usepackage{needspace,float}
\usepackage[font=small,labelfont={sf,color=Bleu},labelsep=period,skip=6pt]{caption}
\usepackage{tikz}
\usetikzlibrary{arrows.meta,calc,decorations.markings}
\usepackage{pgfplots}
\pgfplotsset{compat=1.17}
\usepackage{tkz-tab}
\usepackage[most]{tcolorbox}
\usepackage{titlesec,titletoc}
\usepackage{fancyhdr,lastpage,etoolbox}
\usepackage[colorlinks=true,linkcolor=Bleu,urlcolor=Bleu,citecolor=Bleu,
            pdftitle={Le Collé — Séries entières : sélection de 5 problèmes},pdfauthor={Ahmed Amine Ouqlani}]{hyperref}

\numberwithin{equation}{section}
\numberwithin{figure}{section}
\numberwithin{table}{section}

% ---------- Titres ----------
% Chapitres (sections) : grand numéro, titre en linéale, filet ocre
\titleformat{\section}[display]
  {\normalfont\sffamily\color{Ink}}
  {{\small\color{Gris}\textls[200]{\MakeUppercase{Chapitre}}}\\[-0.2em]%
   {\fontsize{64}{64}\selectfont\color{BleuClair}\thesection}}
  {0.4em}
  {\fontsize{25}{30}\selectfont\bfseries}
  [\vspace{0.5em}{\color{Ocre}\rule{3.2cm}{1.4pt}}]
\titlespacing*{\section}{0pt}{0pt}{1.6em}
\titleformat{name=\section,numberless}[display]
  {\normalfont\sffamily\color{Ink}}{}{0pt}{\fontsize{25}{30}\selectfont\bfseries}
  [\vspace{0.5em}{\color{Ocre}\rule{3.2cm}{1.4pt}}]
\titlespacing*{name=\section,numberless}{0pt}{0pt}{1.6em}

\titleformat{\subsection}
  {\needspace{5\baselineskip}\normalfont\sffamily\Large\bfseries\color{Ink}}
  {\color{Bleu}\thesubsection}{0.8em}{}
\titlespacing*{\subsection}{0pt}{1.8em}{0.5em}
\titleformat{\subsubsection}
  {\needspace{5\baselineskip}\normalfont\sffamily\large\bfseries\color{Ink}}
  {\color{Bleu}\thesubsubsection}{0.7em}{}
\titlespacing*{\subsubsection}{0pt}{1.6em}{0.4em}
\titleformat{\paragraph}[runin]{\normalfont\sffamily\bfseries\color{Ink}}{}{0pt}{}[.]
\titlespacing*{\paragraph}{0pt}{0.9em}{0.6em}
\setcounter{secnumdepth}{3}
\setcounter{tocdepth}{2}

% ---------- Table des matières ----------
\titlecontents{section}[0pt]{\addvspace{0.7em}\sffamily\bfseries\color{Ink}}
  {\makebox[2.2em][l]{\color{Bleu}\thecontentslabel}}{}{\hfill\contentspage}[\addvspace{0.2em}]
\titlecontents{subsection}[2.2em]{\small\linespread{1}\selectfont}
  {\makebox[2.6em][l]{\color{Gris}\thecontentslabel}}{}{\titlerule*[0.7pc]{.}\contentspage}

% ---------- En-têtes et pieds de page ----------
\pagestyle{fancy}
\fancyhf{}
\renewcommand{\sectionmark}[1]{\markboth{\thesection\quad#1}{}}
\fancyhead[L]{\footnotesize\sffamily\color{Gris}\leftmark}
\fancyhead[R]{\footnotesize\sffamily\scshape\color{Gris}Le Collé}
\fancyfoot[C]{\small\sffamily\color{Ink}\thepage}
\fancyfoot[L]{\scriptsize\sffamily\hyperlink{toc}{\color{Filet}sommaire}}
\renewcommand{\headrulewidth}{0.4pt}
\renewcommand{\headrule}{{\color{Filet}\hrule height\headrulewidth width\headwidth}}
\renewcommand{\footrulewidth}{0pt}
\fancypagestyle{plain}{\fancyhf{}\fancyfoot[C]{\small\sffamily\color{Ink}\thepage}\renewcommand{\headrulewidth}{0pt}}

% ---------- Énoncés ----------
\newtheoremstyle{enonce}{0pt}{0pt}{}{}{\sffamily\bfseries\color{Bleu}}{.}{0.5em}
  {\thmname{#1}\thmnumber{~#2}\thmnote{\ \normalfont\sffamily\color{Ink}(#3)}}
\newtheoremstyle{exemple}{0.9em}{0.9em}{}{}{\sffamily\bfseries\color{BleuMoyen}}{.}{0.5em}
  {\thmname{#1}\thmnumber{~#2}\thmnote{\ \normalfont\sffamily\color{Ink}(#3)}}
\newtheoremstyle{remarque}{0.7em}{0.7em}{}{}{\sffamily\itshape}{.}{0.5em}{\thmname{#1}\thmnote{\ (#3)}}
\theoremstyle{enonce}
\newtheorem{definition}{Définition}[section]
\newtheorem{theoreme}[definition]{Théorème}
\newtheorem{proposition}[definition]{Proposition}
\newtheorem{corollaire}[definition]{Corollaire}
\theoremstyle{exemple}
\newtheorem{exemple}[definition]{Exemple}
\theoremstyle{remarque}
\newtheorem*{remarque}{Remarque}

\tcbset{
  defi/.style={enhanced,breakable,colback=white,colframe=white,boxrule=0pt,arc=0pt,
     borderline west={1.2pt}{0pt}{Bleu},left=11pt,right=4pt,top=3pt,bottom=3pt,
     before skip=1.1em,after skip=1.1em},
  thm/.style={enhanced,breakable,colback=Fond,colframe=Fond,boxrule=0pt,arc=0pt,
     borderline west={2.4pt}{0pt}{Bleu},left=11pt,right=8pt,top=6pt,bottom=6pt,
     before skip=1.2em,after skip=1.2em}}
\tcolorboxenvironment{definition}{defi}
\tcolorboxenvironment{theoreme}{thm}
\tcolorboxenvironment{proposition}{thm}
\tcolorboxenvironment{corollaire}{thm}

\renewcommand{\proofname}{Démonstration}
\renewcommand{\qedsymbol}{\textcolor{BleuClair}{$\blacksquare$}}
\makeatletter
\renewenvironment{proof}[1][\proofname]{\par\pushQED{\qed}\normalfont
  \topsep0.5em\relax\trivlist\item[\hskip\labelsep\sffamily\itshape\color{Gris}#1.]\ignorespaces}
  {\popQED\endtrivlist\@endpefalse}
\makeatother

% ---------- Méthode : filets haut et bas, titre en bandeau ----------
\makeatletter\newcommand{\methlabelfix}{\protected@edef\@currentlabel{\thetcbcounter}}\makeatother
\newtcolorbox[auto counter,number within=section]{methode}[2][]{%
  enhanced,unbreakable,colback=white,colframe=white,boxrule=0pt,arc=0pt,
  borderline north={0.9pt}{0pt}{Bleu},borderline south={0.9pt}{0pt}{Bleu},
  left=4pt,right=4pt,top=8pt,bottom=8pt,before skip=1.4em,after skip=1.4em,
  before upper={\methlabelfix\phantomsection\addcontentsline{met}{metentry}{\protect\makebox[2.6em][l]{\thetcbcounter}#2}%
     {\sffamily\bfseries\small\color{white}\colorbox{Bleu}{\strut\ MÉTHODE~\thetcbcounter\ }}%
     \ \ {\sffamily\bfseries\color{Bleu}#2}\par\smallskip},#1}

% ---------- Repères pédagogiques : un filet de couleur, un mot-clé ----------
\newcommand{\definirbloc}[3]{% nom, couleur, libellé
  \newtcolorbox{#1}[1][]{enhanced,breakable,blanker,
    borderline west={1.4pt}{0pt}{#2},left=10pt,right=2pt,top=1pt,bottom=1pt,
    before skip=1em,after skip=1em,
    before upper={{\sffamily\bfseries\footnotesize\color{#2}\textls[120]{\MakeUppercase{#3}}}\quad},##1}}
\definirbloc{reflexe}{Bleu}{Réflexe}
\definirbloc{astuce}{Ocre}{Astuce}
\definirbloc{attention}{Brique}{Attention}
\definirbloc{enpratique}{BleuMoyen}{En pratique}
\definirbloc{plusloin}{Gris}{Pour aller plus loin}
\definirbloc{strategie}{Bleu}{Stratégie}
\definirbloc{lecture}{BleuMoyen}{Lire le schéma}
% compatibilité avec les parties reprises
\newenvironment{garde}{\begin{attention}}{\end{attention}}
\newenvironment{complementHP}{\begin{plusloin}}{\end{plusloin}}
\newenvironment{interpretation}{\begin{enpratique}}{\end{enpratique}}

% ---------- Exercices, corrigés ----------
\newcounter{exo}[section]
\renewcommand{\theexo}{\thesection.\arabic{exo}}
\newcommand{\niveau}[1]{{\color{Ocre}\footnotesize
  \ifcase#1\or$\bullet\circ\circ\circ\circ$\or$\bullet\bullet\circ\circ\circ$\or$\bullet\bullet\bullet\circ\circ$%
  \or$\bullet\bullet\bullet\bullet\circ$\or$\bullet\bullet\bullet\bullet\bullet$\fi}}
\newcommand{\serie}[1]{\par\needspace{8\baselineskip}\vspace{1.5em}%
  \noindent{\sffamily\bfseries\color{Bleu}#1}\par\nobreak\vspace{-0.4em}%
  \noindent{\color{Filet}\rule{\linewidth}{0.4pt}}\par\nobreak\vspace{0.1em}}
\newenvironment{exo}[3]{%
  \par\needspace{5\baselineskip}\vspace{0.75em}\refstepcounter{exo}\label{e:#1}%
  \def\exolab{#1}%
  \noindent{\sffamily\bfseries\color{Bleu}Exercice~\theexo}\quad\niveau{#2}%
  \hfill{\footnotesize\sffamily\color{Gris}#3}\par\nobreak\vspace{0.15em}}%
  {\par\nopagebreak\noindent\hfill{\scriptsize\sffamily\hyperref[c:\exolab]{\color{Gris}corrigé p.~\pageref{c:\exolab}}}\par}
\newenvironment{corrige}[1]{%
  \par\needspace{6\baselineskip}\vspace{1em}\phantomsection\label{c:#1}%
  \noindent{\sffamily\bfseries\color{Bleu}Exercice~\ref{e:#1}}\quad
  {\scriptsize\sffamily\hyperref[e:#1]{\color{Gris}énoncé p.~\pageref{e:#1}}}\par\nobreak\vspace{0.2em}}%
  {\hfill\textcolor{BleuClair}{$\square$}\par}
\newcommand{\etape}[1]{\par\smallskip\noindent{\sffamily\bfseries\small\color{BleuMoyen}#1.}\ }
\newenvironment{solution}{\par\noindent{\sffamily\bfseries\small\color{BleuMoyen}Solution.}\ }{\hfill\textcolor{BleuClair}{$\square$}\par}
\newcommand{\chercher}{\par\smallskip\noindent{\color{Filet}\rule[0.5ex]{1.2cm}{0.4pt}}\
  {\footnotesize\sffamily\itshape\color{Gris}laisser chercher l'étudiant avant de poursuivre}\par\smallskip}

% ---------- Guide de séance ----------

\newcommand{\marge}[1]{\leavevmode\marginpar{\raggedleft\scriptsize\sffamily\color{Gris}\textls[100]{\MakeUppercase{#1}}}\ignorespaces}
\newcommand{\Etape}[4]{%
  \subsection[Étape #1 --- #3 {\color{Gris}(#2)}]{#3}\label{etape:#1}%
  \noindent{\small\sffamily\color{Gris}{\color{Ocre}\bfseries Étape #1}\quad·\quad #2\quad·\quad #4}\par\medskip}
\newcommand{\question}[1]{\par\smallskip\noindent{\sffamily\bfseries\small\color{BleuMoyen}Question.}\ \textit{#1}\par\smallskip}
\newcommand{\rappel}[1]{\par\smallskip\noindent\phantomsection\addcontentsline{rap}{rapentry}{#1}%
  {\sffamily\bfseries\scriptsize\colorbox{Fond}{\color{Gris}\strut RAPPEL}}\ {\sffamily\small\color{Gris}#1.}\ }
\newcommand{\pourquoi}{\par\noindent{\sffamily\bfseries\small\color{Ocre}Pourquoi ?}\ }

% Ouverture de chapitre : \ouverture{objectifs}{prérequis}{durée}
\newcommand{\ouverture}[3]{%
  \begin{tcolorbox}[enhanced,colback=Fond,colframe=Fond,arc=0pt,boxrule=0pt,left=12pt,right=12pt,top=10pt,bottom=10pt,
     before skip=1.2em,after skip=1.6em]
  \begin{minipage}[t]{0.56\linewidth}
  {\sffamily\bfseries\small\color{Bleu}OBJECTIFS}\par\vspace{0.3em}\small #1
  \end{minipage}\hfill
  \begin{minipage}[t]{0.38\linewidth}
  {\sffamily\bfseries\small\color{Bleu}PRÉREQUIS}\par\vspace{0.3em}\small #2
  \ifx\relax#3\relax\else\par\vspace{0.6em}{\sffamily\bfseries\small\color{Bleu}DURÉE}\par\vspace{0.2em}\small #3\fi
  \end{minipage}
  \end{tcolorbox}}
% Fin de chapitre
\newenvironment{essentiel}{\par\needspace{12\baselineskip}\vspace{1.5em}%
  \begin{tcolorbox}[enhanced,unbreakable,colback=white,colframe=white,boxrule=0pt,arc=0pt,
   borderline north={1.4pt}{0pt}{Ocre},left=2pt,right=2pt,top=8pt,bottom=4pt]
  {\sffamily\bfseries\color{Ocre}L'ESSENTIEL DU CHAPITRE}\par\medskip}{\end{tcolorbox}}

% ---------- Notations ----------
\newcommand{\K}{\mathbb{K}}
\newcommand{\R}{\mathbb{R}}
\newcommand{\C}{\mathbb{C}}
\newcommand{\N}{\mathbb{N}}
\newcommand{\Z}{\mathbb{Z}}
\newcommand{\Mn}[1]{\mathcal{M}_{n}(#1)}
\newcommand{\GLn}[1]{\mathrm{GL}_n(#1)}
\newcommand{\LE}{\mathcal{L}(E)}
\newcommand{\rg}{\operatorname{rg}}
\newcommand{\vect}{\operatorname{Vect}}
\newcommand{\Mat}{\operatorname{Mat}}
\newcommand{\Id}{\operatorname{Id}}
\newcommand{\sh}{\operatorname{sh}}
\newcommand{\ch}{\operatorname{ch}}
\renewcommand{\th}{\operatorname{th}}
\DeclareMathOperator{\argsh}{argsh}
\newcommand{\vv}[1]{\overrightarrow{#1}}
\newcommand{\ve}[1]{\vec{#1}}
\newcommand{\ient}[2]{\llbracket #1,#2\rrbracket}
\newcommand{\dt}[1]{\left|\begin{matrix}#1\end{matrix}\right|}
\newcommand{\dis}{\displaystyle}

% ---------- Figures ----------
\pgfplotsset{trig format plots=rad}
\pgfplotsset{courbe/.style={axis lines=middle,axis equal image,
  axis line style={Gris,-{Stealth[length=5pt]}},
  xlabel={$x$},ylabel={$y$},ticklabel style={font=\scriptsize\sffamily,color=Gris},
  every axis x label/.style={at={(ticklabel* cs:1)},anchor=west,font=\small},
  every axis y label/.style={at={(ticklabel* cs:1)},anchor=south,font=\small},
  samples=300,smooth,width=0.6\linewidth}}
\tikzset{arcetude/.style={Bleu,line width=1.25pt},
         reste/.style={BleuClair,line width=1pt},
         tgte/.style={Ocre,line width=0.9pt},
         asy/.style={Gris,dashed,line width=0.6pt},
         pt/.style={circle,fill=Ink,inner sep=1.2pt}}

% ---------- Navigation ----------
\makeatletter
\newcommand*\l@rapentry{\@dottedtocline{1}{0em}{0em}}
\newcommand*\l@metentry{\@dottedtocline{1}{0em}{0em}}
\newcommand{\listemethodes}{\@starttoc{met}}
\newcommand{\listerappels}{\@starttoc{rap}}
\makeatother
% onglets de chapitre sur le bord droit
\usepackage{eso-pic}
\newif\ifonglets \ongletsfalse
\newcommand{\dessineronglet}{%
  \edef\ongcol{\ifnum\value{section}=3 Ocre\else Bleu\fi}%
  \edef\ongnum{\the\value{section}}%
  \pgfmathsetmacro{\ongy}{3.2+2.3*(\ongnum-1)}%
  \begin{tikzpicture}[remember picture,overlay]
  \fill[\ongcol] ([xshift=-0.75cm,yshift=-\ongy cm]current page.north east)
     rectangle ([yshift=-\ongy cm-1.9cm]current page.north east);
  \node[text=white,font=\sffamily\bfseries\large] at ([xshift=-0.37cm,yshift=-\ongy cm-0.95cm]current page.north east)
     {\ongnum};
  \end{tikzpicture}}
\AddToShipoutPictureBG{\ifonglets\ifnum\value{section}>0 \dessineronglet\fi\fi}
% blocs de navigation et d'enseignement
\definirbloc{commencer}{Gris}{Si vous commencez ici}
\definirbloc{horsprogramme}{Brique}{Hors programme}
\newcommand{\HP}{\,{\sffamily\bfseries\scriptsize\textcolor{Brique}{[HP]}}}
\newcommand{\source}[1]{{\footnotesize\sffamily\color{Gris}\,(#1)}}
\newcommand{\coupdepouce}[1]{\par\vspace{-0.2em}\noindent{\footnotesize\sffamily\color{Gris}%
  \textbf{Coup de pouce} (à donner si l'étudiant bloque) : #1}\par}
\newcommand{\pointarret}[1]{\par\vspace{1em}\noindent{\color{Filet}\rule[0.5ex]{\linewidth}{0.4pt}}\par\nobreak
  \noindent{\footnotesize\sffamily\color{Ocre}\bfseries POINT D'ARRÊT POSSIBLE}\quad{\footnotesize\sffamily\color{Gris}#1}\par\vspace{0.6em}}
\newcommand{\demarche}[1]{\noindent{\sffamily\bfseries\small\color{BleuMoyen}Démarche.}\ \textit{#1}\par\smallskip}

% ---------- Corrigés : idées, astuces, pièges, à retenir ----------
\newcommand{\choix}[1]{\par\noindent{\sffamily\bfseries\small\color{BleuMoyen}Choix de la méthode.}\ #1\par\smallskip}
\newcommand{\gain}[1]{\par\noindent{\sffamily\bfseries\small\color{Ocre}Astuce.}\ #1\par}
\newcommand{\piege}[1]{\par\noindent{\sffamily\bfseries\small\color{Brique}Piège.}\ #1\par}
\newtcolorbox{retenirbox}{enhanced,breakable,colback=Fond,colframe=Fond,arc=0pt,boxrule=0pt,
  borderline west={2pt}{0pt}{Ocre},left=8pt,right=6pt,top=4pt,bottom=4pt,before skip=0.7em,after skip=0.5em}
\newcommand{\retenir}[1]{\begin{retenirbox}{\sffamily\bfseries\footnotesize\color{Ocre}À RETENIR}\quad #1\end{retenirbox}}
\newcommand{\TD}[1]{{\color{Ocre}\bfseries TD\,#1}\ ·\ }
\newcommand{\q}[1]{\par\smallskip\noindent{\sffamily\bfseries\color{Bleu}#1}\ }
\providecommand{\iddots}{\mathinner{\mkern1mu\raise1pt\hbox{.}\mkern2mu\raise4pt\hbox{.}\mkern2mu\raise7pt\hbox{.}\mkern1mu}}

\providecommand{\Ker}{\operatorname{Ker}}\renewcommand{\Im}{\operatorname{Im}}\providecommand{\tr}{\operatorname{tr}}\providecommand{\Sp}{\operatorname{Sp}}

% ================================================================ TD : mise en page propre au livret d'exercices
\fancyhf{}
\fancyhead[L]{\footnotesize\sffamily\color{Gris}Sélection · Séries entières}
\fancyhead[R]{\footnotesize\sffamily\scshape\color{Gris}Le Collé}
\fancyfoot[C]{\small\sffamily\color{Ink}\thepage\,/\,\pageref{LastPage}}
\fancypagestyle{plain}{\fancyhf{}\fancyfoot[C]{\small\sffamily\color{Ink}\thepage\,/\,\pageref{LastPage}}\renewcommand{\headrulewidth}{0pt}}
\makeatletter
\newcommand{\partie}[3]{% numéro, titre, description
  \par\needspace{10\baselineskip}\vspace{1.4em}\phantomsection\label{partie:#1}\pdfbookmark[0]{Partie #1}{partie.#1}%
  {\sffamily\small\color{Gris}\textls[200]{PARTIE #1}}\par\nobreak\vspace{0.1em}
  {\sffamily\bfseries\LARGE\color{Ink}#2}\par\nobreak\vspace{0.3em}
  {\color{Ocre}\rule{2.6cm}{1.2pt}}\par\nobreak\vspace{0.2em}
  {\small\color{Gris}#3}\par\nobreak\vspace{0.6em}}
\newcommand{\exercice}[7]{% numéro, difficulté, titre, origine, notions, identifiant, page du corrigé
  \par\needspace{8\baselineskip}\vspace{1.1em}\phantomsection
  \def\@currentlabel{#1}\label{e:#6}\pdfbookmark[1]{Exercice #1}{e.#6}%
  \noindent{\sffamily\bfseries\large\color{Bleu}Exercice~#1}\quad\niveau{#2}\hfill{\footnotesize\sffamily\color{Gris}#4}\par\nobreak
  \vspace{0.1em}\noindent{\sffamily\bfseries\color{Ink}#3}\par\nobreak
  \vspace{0.1em}\noindent{\footnotesize\sffamily\color{Gris}Notions : #5.\hfill\mbox{corrigé p.~#7}\par}\nobreak
  \vspace{0.35em}}
\newcommand{\renvoi}[1]{\par\nopagebreak\vspace{0.2em}\noindent\hfill{\scriptsize\sffamily\color{Gris}#1}\par}
\newcommand{\place}{\par\vspace{1.6cm plus 0.6cm minus 0.6cm}}
\newcommand{\correction}[4]{% numéro, titre, identifiant, renvoi
  \par\needspace{7\baselineskip}\vspace{1.3em}\phantomsection
  \def\@currentlabel{#1}\label{e:#3}\pdfbookmark[1]{Exercice #1}{c.#3}%
  {\color{Filet}\hrule height 0.4pt}\vspace{0.5em}
  \noindent{\sffamily\bfseries\large\color{Bleu}Exercice~#1}\quad{\sffamily\bfseries\color{Ink}#2}\hfill{\scriptsize\sffamily\color{Gris}#4}\par\nobreak
  \vspace{0.4em}}
\newcommand{\sommaireligne}[3]{\noindent\hbox to 2.2em{\hfill\sffamily\color{Bleu}#1.}\enspace\begin{minipage}[t]{0.8\linewidth}\raggedright#2\end{minipage}\hfill\pageref{e:#3}\par}
\makeatother
\setlength{\parskip}{0.5em plus 0.15em minus 0.1em}

\makeatletter
\newcommand{\probleme}[7]{% numéro, étape, titre, difficulté, origine, durée, identifiant
  \clearpage\phantomsection\label{p:#7}\pdfbookmark[1]{Problème #1 — #2}{p.#7}%
  {\sffamily\small\color{Gris}\textls[150]{PROBLÈME #1\ \ ·\ \ \MakeUppercase{#2}}}\par\nobreak\vspace{0.15em}
  {\sffamily\bfseries\LARGE\color{Ink}#3}\par\nobreak\vspace{0.3em}
  {\color{Ocre}\rule{2.6cm}{1.2pt}}\par\nobreak\vspace{0.2em}
  \noindent{\footnotesize\sffamily\color{Gris}Difficulté~\niveau{#4}\quad·\quad durée indicative : #6\quad·\quad #5}\par\nobreak\vspace{0.3em}}
\newcommand{\fichepb}[3]{% notions, objectifs, prérequis
  {\small\setlength{\parskip}{0.15em}%
  \noindent\hangindent=2.4cm\hangafter=1\makebox[2.4cm][l]{\sffamily\textcolor{Bleu}{Notions}}#1\par
  \noindent\hangindent=2.4cm\hangafter=1\makebox[2.4cm][l]{\sffamily\textcolor{Bleu}{Objectifs}}#2\par
  \noindent\hangindent=2.4cm\hangafter=1\makebox[2.4cm][l]{\sffamily\textcolor{Bleu}{Prérequis}}#3\par}
  \vspace{0.4em}{\color{Filet}\hrule height 0.4pt}\vspace{0.6em}}
\newcommand{\titresection}[2]{% titre, texte
  \clearpage\phantomsection\pdfbookmark[0]{#1}{s.#1}%
  {\sffamily\bfseries\fontsize{24}{28}\selectfont\color{Ink}#1}\par\vspace{0.3em}
  {\color{Ocre}\rule{3.2cm}{1.4pt}}\par\vspace{0.5em}{\small\color{Gris}#2}\par\vspace{0.6em}}
\newcommand{\soustitre}[3]{% numéro, titre, label
  \par\needspace{6\baselineskip}\vspace{1em}\phantomsection\label{#3}%
  {\color{Filet}\hrule height 0.4pt}\vspace{0.5em}
  \noindent{\sffamily\bfseries\large\color{Bleu}Problème~#1}\quad{\sffamily\bfseries\color{Ink}#2}\par\nobreak\vspace{0.3em}}
\makeatother

\begin{document}
\thispagestyle{plain}
{\sffamily\small\color{Gris}\textls[150]{LE COLLÉ · MATHÉMATIQUES PT · CHAPITRE 8}}\par\vspace{0.2em}
{\sffamily\bfseries\fontsize{26}{30}\selectfont\color{Ink}Sélection --- Séries entières}\par\vspace{0.3em}
{\sffamily\Large\color{BleuMoyen}Cinq problèmes à travailler après le cours}\par\vspace{0.4em}
{\color{Ocre}\rule{3.2cm}{1.4pt}}\par\vspace{0.8em}
{\small Ces cinq problèmes sont la sélection principale pour les séances, \emph{après} l'explication du cours.
Ils vont de la mise en application au niveau des devoirs et des concours ; chacun réunit plusieurs résultats du chapitre.
Ils complètent le TD complet et le livre, qu'ils ne remplacent pas.

\textbf{Mode d'emploi.} Les traiter dans l'ordre. Pour chaque problème : lire l'en-tête (notions, objectifs, prérequis),
chercher au brouillon, puis rédiger. En cas de blocage, lire \emph{une seule} indication (page des indications), puis reprendre.
Ne lire la correction qu'une fois le problème terminé ou après un vrai effort : elle explique le \emph{choix de la méthode}
avant les calculs. L'origine de chaque problème est indiquée en tête ; les problèmes « d'après » un concours sont des
adaptations rédigées par Le Collé, pas les énoncés officiels.}
\vspace{1.2em}
{\small\renewcommand{\arraystretch}{1.45}\setlength{\tabcolsep}{4pt}\noindent\begin{tabular}{@{}r>{\raggedright\arraybackslash\sffamily\footnotesize}p{2.9cm}>{\raggedright\arraybackslash}p{6.1cm}ll@{\hspace{6pt}}ccc@{}}
& & & & & {\footnotesize\sffamily énoncé} & {\footnotesize\sffamily indic.} & {\footnotesize\sffamily corrigé}\\\hline
1 & \textcolor{Gris}{Mise en application} & Rayons de convergence & \niveau{2} & 25\,min & \pageref{p:se-rayons} & \pageref{i:se-rayons} & \pageref{c:se-rayons}\\
2 & \textcolor{Gris}{Méthodes classiques} & Continuité de la somme au bord & \niveau{2} & 30\,min & \pageref{p:set-12} & \pageref{i:set-12} & \pageref{c:set-12}\\
3 & \textcolor{Gris}{Approfondissement} & Développements par une équation différentielle & \niveau{3} & 40\,min & \pageref{p:set-16} & \pageref{i:set-16} & \pageref{c:set-16}\\
4 & \textcolor{Gris}{Problème de synthèse} & Une suite récurrente et sa série entière & \niveau{4} & 45\,min & \pageref{p:set-21} & \pageref{i:set-21} & \pageref{c:set-21}\\
5 & \textcolor{Gris}{Niveau DS / concours} & Équation différentielle du second ordre & \niveau{4} & 45\,min & \pageref{p:set-22} & \pageref{i:set-22} & \pageref{c:set-22}\\
\end{tabular}}
\probleme{1}{Mise en application}{Rayons de convergence}{2}{Livret, exercice 1.5 · TD complet, exercice~9}{25\,min}{se-rayons}
\fichepb{comparaison, équivalents, d'Alembert, série lacunaire.}{Choisir, d'après la forme du coefficient, l'outil qui donne le rayon : encadrement, équivalent ou règle de d'Alembert.}{Définition et caractérisation du rayon ; Règle de d'Alembert pour les séries numériques ; Rayon de $\sum n^\alpha z^n$.}
Déterminer le rayon de convergence des séries entières :
\begin{gather*}\text{(a)}\ \sum e^{\cos n}z^n\qquad\text{(b)}\ \sum\frac{e^{in}}{1+in}z^n\qquad\text{(c)}\ \sum_{n\geq2}\frac{(-1)^n}{n(n-1)}z^n\\
\text{(d)}\ \sum\frac{(2n)!}{(n!)^2}z^n\qquad\text{(e)}\ \sum2^nz^{3n}.\end{gather*}
\par\vspace{0.8em}\noindent\hfill{\scriptsize\sffamily\color{Gris}indications p.~\pageref{i:se-rayons} · correction p.~\pageref{c:se-rayons}}
\probleme{2}{Méthodes classiques}{Continuité de la somme au bord}{2}{TD de la classe, exercice 12 · TD complet, exercice~15}{30\,min}{set-12}
\fichepb{éléments simples, logarithme, continuité au bord, série télescopique.}{Calculer une fonction somme sur l'intervalle ouvert, puis en déduire des sommes numériques par passage à la limite au bord.}{DSE de $\ln(1+x)$ ; Continuité de la somme au bord (admise) ; Décomposition en éléments simples.}
On considère $\sum_{n\geq2}\frac{(-1)^n}{n(n-1)}x^n$ et on note $f$ sa fonction somme.
\begin{enumerate}
\item Déterminer le rayon de convergence $R$ et l'ensemble de définition de $f$.
\item Calculer $f(x)$ pour tout $x\in\,]-R,R[$.
\item En déduire la nature et la somme des séries numériques $\sum_{n\geq2}\frac1{n(n-1)}$ et $\sum_{n\geq2}\frac{(-1)^n}{n(n-1)}$.
\end{enumerate}
\par\vspace{0.8em}\noindent\hfill{\scriptsize\sffamily\color{Gris}indications p.~\pageref{i:set-12} · correction p.~\pageref{c:set-12}}
\probleme{3}{Approfondissement}{Développements par une équation différentielle}{3}{TD de la classe, exercice 16 · TD complet, exercice~25}{40\,min}{set-16}
\fichepb{équation différentielle, relation de récurrence, unicité du DSE, unicité de Cauchy.}{Mettre en œuvre complètement la méthode analyse-synthèse : équation vérifiée par $f$, récurrence, rayon, conclusion par Cauchy.}{Dérivation terme à terme ; Théorème de Cauchy linéaire ; Dérivées de $\arcsin$.}
En utilisant la technique de l'équation différentielle, trouver les développements en série entière des fonctions suivantes.
\begin{enumerate}
\item $f(x)=\frac{\arcsin x}{\sqrt{1-x^2}}$.
\item $f(x)=\arcsin^2x$.
\item $f(x)=\cos(x)\ch(x)$ : vérifier que $f$ est solution de l'équation $y^{(4)}+4y=0$.
\item $f(x)=\cos(\alpha\arcsin x)$, $\alpha\in\R$ : vérifier que $f$ est solution d'une équation différentielle d'ordre 2.
\end{enumerate}
\par\vspace{0.8em}\noindent\hfill{\scriptsize\sffamily\color{Gris}indications p.~\pageref{i:set-16} · correction p.~\pageref{c:set-16}}
\probleme{4}{Problème de synthèse}{Une suite récurrente et sa série entière}{4}{TD de la classe, exercice 21 · TD complet, exercice~32}{45\,min}{set-21}
\fichepb{récurrence double, rayon par comparaison, équation différentielle du premier ordre, intégration terme à terme.}{Transformer une récurrence sur les coefficients en une équation fonctionnelle, puis différentielle, et obtenir la somme.}{Comparaison des rayons ; Résolution du premier ordre ; Intégration terme à terme.}
Soit $(a_n)$ définie par $a_0=a_1=1$ et $a_{n+2}=a_{n+1}+\frac{a_n}{n+2}$ pour $n\geq0$. On considère $\sum_{n\geq0}a_nx^n$.
\begin{enumerate}
\item (a) Montrer que pour tout $n$, $1\leq a_n\leq n+2$.\quad (b) En déduire que $(a_n)$ est monotone et déterminer le rayon de convergence.
\item On pose $f(x)=\sum_{n=0}^{+\infty}a_nx^n$ et $g(x)=\sum_{n\geq0}\frac{a_n}{n+2}x^{n+2}$.
\begin{enumerate}
\item Justifier que $f$ et $g$ sont de classe $\mathcal C^\infty$ sur $]-1,1[$.
\item Trouver une relation entre $f$ et $g$.
\item Montrer que $g$ est solution d'une équation différentielle linéaire d'ordre 1 que l'on résoudra explicitement.
\item Donner une expression de $f(x)$.
\end{enumerate}
\end{enumerate}
\par\vspace{0.8em}\noindent\hfill{\scriptsize\sffamily\color{Gris}indications p.~\pageref{i:set-21} · correction p.~\pageref{c:set-21}}
\probleme{5}{Niveau DS / concours}{Équation différentielle du second ordre}{4}{TD de la classe, exercice 22 · TD complet, exercice~33}{45\,min}{set-22}
\fichepb{série entière solution, récurrence, rayon, formule de Stirling.}{Chercher une solution développable d'une équation dont le coefficient dominant s'annule, puis étudier ses coefficients : le type de raisonnement des épreuves écrites.}{Unicité du DSE ; Règle de d'Alembert ; Formule de Stirling.}
On considère l'équation $16(x^2-x)y''+(16x-8)y'-y=0$ $(E)$. On suppose que la série entière $\sum a_nx^n$ est solution de $(E)$ sur $]0,R[$, où
$R>0$ est son rayon de convergence.
\begin{enumerate}
\item Montrer que pour tout $n\geq1$, $a_n=\frac{(4n-3)(4n-5)}{8n(2n-1)}a_{n-1}$.
\item Déterminer le rayon de convergence de la série entière $\sum a_nx^n$.
\item On suppose que $a_0=1$. Déterminer $a_n$ en fonction de $n$.
\item En utilisant la formule de Stirling, déterminer un équivalent de $a_n$ lorsque $n\to+\infty$.
\end{enumerate}
\par\vspace{0.8em}\noindent\hfill{\scriptsize\sffamily\color{Gris}indications p.~\pageref{i:set-22} · correction p.~\pageref{c:set-22}}
\titresection{Indications}{Une indication par question délicate. N\'en lire qu\'une à la fois, puis reprendre la recherche.}
\soustitre{1}{Rayons de convergence}{i:se-rayons}
\begin{itemize}[leftmargin=1.2em,itemsep=0.25em]\item (a) Encadrer $e^{\cos n}$ entre deux constantes.
\item (b) Calculer le module $\left|\frac{e^{in}}{1+in}\right|$.
\item (c) Un équivalent du coefficient, puis la proposition des équivalents.
\item (d) Règle de d'Alembert : $\frac{a_{n+1}}{a_n}$ se simplifie.
\item (e) Quotient de deux termes \emph{présents} : $\frac{|2^{n+1}z^{3n+3}|}{|2^nz^{3n}|}$.
\end{itemize}
\soustitre{2}{Continuité de la somme au bord}{i:set-12}
\begin{itemize}[leftmargin=1.2em,itemsep=0.25em]\item 1. $|a_n|\sim\frac1{n^2}$ ; étudier la convergence absolue en $\pm1$.
\item 2. $\frac1{n(n-1)}=\frac1{n-1}-\frac1n$ ; reconnaître deux fois $\ln(1+x)$.
\item 3. Les deux séries numériques sont $f(-1)$ et $f(1)$.
\end{itemize}
\soustitre{3}{Développements par une équation différentielle}{i:set-16}
\begin{itemize}[leftmargin=1.2em,itemsep=0.25em]\item 1. Montrer que $(1-x^2)f'-xf=1$, puis identifier le coefficient de $x^n$.
\item 2. $\arcsin^2$ a pour dérivée $2\frac{\arcsin x}{\sqrt{1-x^2}}$ : intégrer le 1.
\item 3. Dériver quatre fois $\cos x\ch x$ ; récurrence de pas $4$.
\item 4. Montrer $(1-x^2)f''-xf'+\alpha^2f=0$ ; récurrence de pas $2$.
\end{itemize}
\soustitre{4}{Une suite récurrente et sa série entière}{i:set-21}
\begin{itemize}[leftmargin=1.2em,itemsep=0.25em]\item 1.(a) Récurrence double, en supposant les deux encadrements aux rangs $n$ et $n+1$.
\item 2.(b) Dériver $g$ terme à terme.
\item 2.(c) Multiplier la relation par $x^{n+2}$ et sommer : $(1-x)f=1+g$.
\item 2.(d) $\frac x{1-x}=-1+\frac1{1-x}$ ; utiliser $g(0)=0$.
\end{itemize}
\soustitre{5}{Équation différentielle du second ordre}{i:set-22}
\begin{itemize}[leftmargin=1.2em,itemsep=0.25em]\item 1. Écrire chaque terme avec $x^n$ ; regrouper $16n^2-1$.
\item 2. Le quotient $\frac{a_n}{a_{n-1}}$ tend vers $1$.
\item 3. Les facteurs $(4k-3)(4k-5)$ parcourent les entiers impairs ; faire apparaître $\binom{4n}{2n}$.
\item 4. $\binom{4n}{2n}\sim\frac{16^n}{\sqrt{2\pi n}}$.
\end{itemize}
\titresection{Corrections}{Chaque correction commence par le choix de la méthode ; les encadrés résument ce qui resservira. Les renvois « méthode~1.4 du livret » désignent le livret du chapitre.}
\soustitre{1}{Rayons de convergence}{c:se-rayons}
\q{(a)} $e^{-1}\leq e^{\cos n}\leq e$ : par comparaison avec $\sum e^{-1}z^n$ et $\sum ez^n$ (rayon $1$), $1\leq R\leq1$. Donc $R=1$.
\q{(b)} $\left|\frac{e^{in}}{1+in}\right|=\frac1{\sqrt{1+n^2}}\sim\frac1n$ : même rayon que $\sum\frac{z^n}n$, $R=1$.
\q{(c)} $\left|\frac{(-1)^n}{n(n-1)}\right|\sim\frac1{n^2}$ : même rayon que $\sum\frac{z^n}{n^2}$, c'est-à-dire $1$ (proposition~1.21 du livret).
\q{(d)} $\frac{a_{n+1}}{a_n}=\frac{(2n+2)(2n+1)}{(n+1)^2}=\frac{2(2n+1)}{n+1}\to4$ : $R=\frac14$.
\q{(e)} Série lacunaire : pour $z\neq0$, $\frac{|2^{n+1}z^{3n+3}|}{|2^nz^{3n}|}=2|z|^3$. Convergence absolue si $2|z|^3<1$, divergence si
$2|z|^3>1$ : $R=2^{-1/3}$.
\par\hfill\textcolor{BleuClair}{$\square$}\par
\soustitre{2}{Continuité de la somme au bord}{c:set-12}
\q{1.} $|a_n|\sim\frac1{n^2}$ : $R=1$. En $x=\pm1$, $\sum\frac1{n(n-1)}$ converge absolument : $f$ est définie sur $[-1,1]$.
\q{2.} $\frac1{n(n-1)}=\frac1{n-1}-\frac1n$. Pour $|x|<1$ : $\sum_{n\geq2}\frac{(-1)^nx^n}{n-1}=x\sum_{k\geq1}\frac{(-1)^{k+1}x^k}k=x\ln(1+x)$ et
$\sum_{n\geq2}\frac{(-1)^nx^n}n=-\ln(1+x)+x$. Donc $f(x)=(1+x)\ln(1+x)-x$.
\q{3.} $f$ est continue sur $[-1,1]$ (proposition~1.28 du livret, admise) :
$\sum_{n\geq2}\frac1{n(n-1)}=f(-1)=\lim_{x\to-1}\big((1+x)\ln(1+x)-x\big)=1$ (on retrouve le télescopage), et $\sum_{n\geq2}\frac{(-1)^n}{n(n-1)}=f(1)=2\ln2-1$.
\par\hfill\textcolor{BleuClair}{$\square$}\par
\soustitre{3}{Développements par une équation différentielle}{c:set-16}
\q{1.} $f'(x)=\frac1{1-x^2}+\frac{x\arcsin x}{(1-x^2)^{3/2}}$, donc $(1-x^2)f'-xf=1$, $f(0)=0$. Avec $f=\sum a_nx^n$, le coefficient de $x^n$ donne
$(n+1)a_{n+1}-na_{n-1}=0$ pour $n\geq1$ et $a_1=1$. Les $a_{2p}$ sont nuls ($a_0=0$) et $a_{2p+1}=\frac{2p}{2p+1}a_{2p-1}$, d'où
$a_{2p+1}=\frac{(2\cdot4\cdots2p)^2}{(2p+1)!}=\frac{4^p(p!)^2}{(2p+1)!}$. Synthèse : $\frac{a_{2p+3}}{a_{2p+1}}x^2\to x^2$, $R=1$ ; sur $]-1,1[$ la somme est solution
du même problème de Cauchy ($1-x^2\neq0$) : $f(x)=\sum_{p\geq0}\frac{4^p(p!)^2}{(2p+1)!}x^{2p+1}$.
\q{2.} $(\arcsin^2)'=2\frac{\arcsin x}{\sqrt{1-x^2}}$ ; en intégrant le 1 (valeur $0$ en $0$) :
$\arcsin^2x=\sum_{p\geq0}\frac{2\cdot4^p(p!)^2}{(2p+2)!}x^{2p+2}=x^2+\frac{x^4}3+\frac{8x^6}{45}+\cdots$ ; $R=1$.
\q{3.} $f'=-\sin x\ch x+\cos x\sh x$, $f''=-2\sin x\sh x$, $f'''=-2(\cos x\sh x+\sin x\ch x)$,
$f^{(4)}=-4\cos x\ch x$. Donc $f^{(4)}+4f=0$, avec $f(0)=1$, $f'(0)=f''(0)=f'''(0)=0$. La récurrence est $a_{n+4}=\frac{-4a_n}{(n+1)(n+2)(n+3)(n+4)}$ :
$a_{4p}=\frac{(-4)^p}{(4p)!}$, les autres sont nuls ; $R=+\infty$ ; unicité de Cauchy : $\cos x\ch x=\sum\frac{(-4)^px^{4p}}{(4p)!}$.
\q{4.} $f'(x)=-\frac{\alpha\sin(\alpha\arcsin x)}{\sqrt{1-x^2}}$, puis $(1-x^2)f''-xf'+\alpha^2f=0$, avec $f(0)=1$, $f'(0)=0$. Coefficient de $x^n$ :
$(n+2)(n+1)a_{n+2}=(n^2-\alpha^2)a_n$. Donc $a_{2p+1}=0$ et $a_{2p}=\frac{\prod_{k=0}^{p-1}\big((2k)^2-\alpha^2\big)}{(2p)!}$. Si $\alpha$ est un entier pair, la somme
est un polynôme ; sinon $\frac{a_{2p+2}}{a_{2p}}x^2\to x^2$ et $R=1$. Sur $]-1,1[$, unicité de Cauchy : $\cos(\alpha\arcsin x)=\sum_{p\geq0}a_{2p}x^{2p}$.
\par\hfill\textcolor{BleuClair}{$\square$}\par
\soustitre{4}{Une suite récurrente et sa série entière}{c:set-21}
\q{1.(a)} Récurrence double. Si $a_n,a_{n+1}\geq1$, alors $a_{n+2}\geq a_{n+1}\geq1$. Si $a_n\leq n+2$ et $a_{n+1}\leq n+3$, alors
$a_{n+2}\leq n+3+\frac{n+2}{n+2}=n+4$.
\q{1.(b)} $a_{n+2}-a_{n+1}=\frac{a_n}{n+2}>0$ : la suite est croissante. Comme $1\leq a_n\leq n+2$, comparaison avec $\sum x^n$ et $\sum(n+2)x^n$
(rayons $1$) : $R=1$.
\q{2.(a)} Ce sont des sommes de séries entières de rayon $1$ ($g$ : intégration terme à terme), donc $\mathcal C^\infty$ sur $]-1,1[$.
\q{2.(b)} $g'(x)=\sum a_nx^{n+1}=xf(x)$.
\q{2.(c)} On multiplie la relation par $x^{n+2}$ et on somme pour $n\geq0$ :
$f(x)-1-x=x\big(f(x)-1\big)+g(x)$, donc $(1-x)f(x)=1+g(x)$. Avec $g'=xf$ : $g'(x)=\frac x{1-x}\big(1+g(x)\big)$, équation linéaire d'ordre 1.
$\frac x{1-x}=-1+\frac1{1-x}$ a pour primitive $-x-\ln(1-x)$ : les solutions sont $g(x)=-1+C\frac{e^{-x}}{1-x}$. Avec $g(0)=0$ : $C=1$,
$g(x)=\frac{e^{-x}}{1-x}-1$.
\q{2.(d)} $f(x)=\frac{1+g(x)}{1-x}=\frac{e^{-x}}{(1-x)^2}$.
\par\hfill\textcolor{BleuClair}{$\square$}\par
\soustitre{5}{Équation différentielle du second ordre}{c:set-22}
\q{1.} On injecte : $16x^2y''=\sum16n(n-1)a_nx^n$, $-16xy''=-\sum16(n+1)na_{n+1}x^n$, $16xy'=\sum16na_nx^n$, $-8y'=-\sum8(n+1)a_{n+1}x^n$.
Le coefficient de $x^n$ est $(16n^2-1)a_n-8(n+1)(2n+1)a_{n+1}$ ; par unicité du DSE (sur $]0,R[$ la somme est nulle, donc tous ses
coefficients aussi), $a_{n+1}=\frac{(4n-1)(4n+1)}{8(n+1)(2n+1)}a_n$, ce qui est la relation demandée au rang $n+1$.
\q{2.} Si $a_0\neq0$, aucun $a_n$ n'est nul et $\frac{a_n}{a_{n-1}}\to\frac{16}{16}=1$ : $R=1$ (si $a_0=0$, la solution est nulle).
\q{3.} Le produit des $(4k-3)(4k-5)$ pour $k=1..n$ vaut $-1\times1\times3\times\cdots\times(4n-3)=-\frac{(4n-2)!}{2^{2n-1}(2n-1)!}$, et celui des
$8k(2k-1)$ vaut $8^nn!\cdot\frac{(2n)!}{2^nn!}=4^n(2n)!$. Donc $a_n=-\frac{2\,(4n-2)!}{16^n(2n-1)!\,(2n)!}=-\frac1{16^n(4n-1)}\binom{4n}{2n}$
(contrôle : $a_1=-\frac18$).
\q{4.} Stirling : $\binom{4n}{2n}\sim\frac{16^n}{\sqrt{2\pi n}}$. Donc $a_n\sim-\frac1{4n\sqrt{2\pi n}}=-\frac1{4\sqrt{2\pi}\,n^{3/2}}$.
\par\hfill\textcolor{BleuClair}{$\square$}\par
\end{document}
